\subsection{Parallel Capture}
\label{sec:parallel-capture}

Parallel capture allows the camera application to accept the next shutter input before the preceding capture has produced its draft image.
This is possible because the hardware abstraction layer (HAL) can issue a capture-availability signal, \texttt{captureAvailable}, before the draft sequence completes.
This signal indicates when the application may accept the next capture request, allowing a shorter shot-to-shot interval.

However, draft sequences execute one at a time on a single draft worker to preserve capture order and limit peak memory demand.
When capture requests arrive faster than the worker completes draft sequences, later captures must wait in the draft queue.
Since each capture's timeout clock starts when its request is accepted, this wait consumes part of its timeout budget before its draft sequence starts.
Figure~\ref{fig:overlap} illustrates how this waiting time grows across successive captures even when their draft sequences take equally long.
Accumulated waiting time can therefore cause later captures to miss their deadlines, especially when lightweight multi-frame composition increases the draft processing time of earlier captures.

\begin{figure}[t]
  \centering
  \resizebox{\columnwidth}{!}{% Notes: docs/exhibits.md#fig_parallel_capture_overlap
\begin{tikzpicture}[font=\small,
  collectstage/.style={draw=black!55, fill=gray!14},
  draftstage/.style={draw=blue!45!black, fill=blue!14},
  queued/.style={pattern=north east lines, pattern color=red!45, draw=red!35},
  overrun/.style={draw=red!80!black, fill=red!65},
  ddl/.style={red, dashed, thick},
  sig/.style={->, black!75}]

  \node[anchor=east] at (-0.2, 3.925) {capture $i$};
  \draw[very thick] (0, 3.55) -- (0, 4.30);
  \draw[collectstage] (0.05, 3.70) rectangle (0.90, 4.15);
  \node at (0.475, 3.925) {\scriptsize collect};
  \draw[draftstage] (0.90, 3.70) rectangle (3.50, 4.15);
  \node at (2.20, 3.925) {\scriptsize Draft Sequence};

  \draw[densely dotted, gray] (3.50, 3.925) -- (7.30, 3.925);
  \node[fill=white, inner sep=1.5pt] at (5.40, 3.925)
    {\footnotesize\textcolor{green!45!black}{\ding{51} in time}};
  \draw[ddl] (7.30, 3.55) -- (7.30, 4.30);
  \node[red, anchor=east] at (9.00, 3.925) {deadline $i$};

  \fill[black!60] (0.35, 3.68) circle (1.3pt);
  \draw[sig] (0.35, 3.66) -- (0.50, 3.17);
  \node[anchor=west, black!75] at (0.42, 3.42)
    {\scriptsize\texttt{captureAvailable}};

  \node[anchor=east] at (0.3, 2.725) {capture $i{+}1$};
  \draw[very thick] (0.50, 2.35) -- (0.50, 3.10);
  \draw[collectstage] (0.55, 2.50) rectangle (1.40, 2.95);
  \node at (0.975, 2.725) {\scriptsize collect};
  \draw[queued] (1.40, 2.50) rectangle (3.50, 2.95);
  \node[fill=white, inner sep=1.5pt] at (2.45, 2.725)
    {\footnotesize\textcolor{red!60!black}{waiting}};
  \draw[draftstage] (3.50, 2.50) rectangle (6.10, 2.95);
  \node at (4.80, 2.725) {\scriptsize Draft Sequence};

  \draw[densely dotted, gray] (6.10, 2.725) -- (7.80, 2.725);
  \node[fill=white, inner sep=1.5pt] at (6.95, 2.725)
    {\footnotesize\textcolor{green!45!black}{\ding{51} in time}};
  \draw[ddl] (7.80, 2.35) -- (7.80, 3.10);
  \node[red, anchor=east] at (10.00, 2.725) {deadline $i+1$};

  \fill[black!60] (0.85, 2.48) circle (1.3pt);
  \draw[sig] (0.85, 2.46) -- (1.00, 1.97);

  \node[anchor=east] at (0.8, 1.525) {capture $i{+}2$};
  \draw[very thick] (1.00, 1.15) -- (1.00, 1.90);
  \draw[collectstage] (1.05, 1.30) rectangle (1.90, 1.75);
  \node at (1.475, 1.525) {\scriptsize collect};
  \draw[queued] (1.90, 1.30) rectangle (6.10, 1.75);
  \node[fill=white, inner sep=1.5pt] at (4.00, 1.525)
    {\footnotesize\textcolor{red!60!black}{waiting}};
  \draw[draftstage] (6.10, 1.30) rectangle (8.70, 1.75);

  \draw[overrun] (8.30, 1.30) rectangle (8.70, 1.75);
  \node at (7.20, 1.525) {\scriptsize Draft Sequence};
  \draw[ddl] (8.30, 1.15) -- (8.30, 1.90);
  \node[red!70!black, anchor=west, align=left] at (8.80, 1.525)
    {\footnotesize\bfseries Capture\\[-2pt]\footnotesize\bfseries Timeout};

  \draw[dashed, gray, ->] (3.50, 3.67) -- (3.50, 2.98);
  \draw[dashed, gray, ->] (6.10, 2.47) -- (6.10, 1.78);

  \draw[decorate, decoration={brace, mirror, amplitude=4pt}]
    (1.00, 1.10) -- (6.10, 1.10);
  \node[anchor=north] at (3.55, 0.88) {\footnotesize Timeout budget consumed};
  \draw[decorate, decoration={brace, mirror, amplitude=4pt}]
    (6.10, 1.10) -- (8.30, 1.10);
  \node[anchor=north] at (7.20, 0.88) {\footnotesize remaining budget};

  \draw[->] (-1.3, 0.42) -- (9.7, 0.42) node[anchor=south east] {time};

\end{tikzpicture}
}
  \caption{Timeout budget consumed by queueing under parallel capture.}
  \label{fig:overlap}
\end{figure}
